\section{Fields of classes and bounded-degree finiteness}\label{sec:fields}

\subsection{Descent by trace}\label{sec:class-fields}

Let \(\mathcal G\) be a rational geometric Chow quotient with its
discrete \(\Gal(\Fbar/F_0)\)-action, and let \(\Lambda(E)\subset\mathcal G\)
be compatible integral images. We assume that restriction and trace
preserve these images in their usual directions. Fixed rational Chow
projectors and quotient by geometric algebraic equivalence have this
property. Every element considered is defined over a finite extension,
so its stabilizer is open. Define its minimal class field by
\[
 F_0(\xi)=\Fbar^{\operatorname{Stab}_{\Gal(\Fbar/F_0)}(\xi)}.
\]

\begin{lemma}\label[lemma]{lem:class-trace}
If \(\xi\in\Lambda(E)\) and \(K=F_0(\xi)\), then
\[
 K\subseteq E,\qquad [E:K]\xi\in\Lambda(K).
\]
In particular, if \([E:F_0]\le D\), the required multiplier is
at most \(D\).
\end{lemma}

\begin{proof}
The inclusion \(\Gal(\Fbar/E)\subseteq\operatorname{Stab}(\xi)\) gives
\(K\subseteq E\). Take an integral cycle over \(E\) representing
\(\xi\) in the indicated quotient and push it forward to \(K\).
Over \(\Fbar\), the resulting class is the sum over the
\(K\)-embeddings of \(E\). Every summand equals \(\xi\), so
the sum is \([E:K]\xi\). Fixed projectors commute with trace,
and algebraic equivalence is preserved by it.
\end{proof}

The lemma concerns the field of a class in a rational quotient.
It supplies the precise denominator even when no chosen integral
cycle representative descends to that field.

\subsection{The ramification criterion}\label{sec:ramification-proof}

Let \(q\) be a Galois-invariant quadratic form on \(\mathcal G\).
For \(B\ge0\) and \(D\ge1\), consider
\[
 \mathcal B(B,D)=
 \{\xi\in\bigcup_{[E:F_0]\le D}\Lambda(E):q(\xi)\le B\}.
\]

\begin{theorem}\label[theorem]{thm:ramification}
Assume \(q\) has fixed-field Northcott on \(\Lambda(K)\) for
every finite \(K/F_0\). Then bounded-degree Northcott is
equivalent to the following condition: for every \(B,D\), there
is a finite set \(S(B,D)\) of finite places of \(F_0\) such that
each field \(F_0(\xi)/F_0\), with \(\xi\in\mathcal B(B,D)\),
is unramified outside \(S(B,D)\).
\end{theorem}

\begin{proof}
Suppose the ramification condition holds. By the finite-ramification
form of Hermite--Minkowski, there are only finitely many extensions
\(K/F_0\) in \(\Fbar\) of degree at most \(D\) unramified
 outside \(S(B,D)\); see \cite[Theorem~10.2.2]{Kedlaya}.
This form also follows by bounding the local discriminant exponents
in fixed degree and applying discriminant finiteness.

For \(\xi\in\mathcal B(B,D)\), \Cref{lem:class-trace} gives
\(m\xi\in\Lambda(K)\), where \(K=F_0(\xi)\) is among these
finitely many fields and \(1\le m\le D\). Moreover
\[
 q(m\xi)=m^2q(\xi)\le D^2B.
\]
Fixed-field Northcott leaves finitely many possibilities for
\(m\xi\). Multiplication by \(m\) is injective in the rational
group, and only finitely many \(m\) occur. Thus the ball is finite.
Conversely, if the ball is finite, the union of the ramification
sets of its finitely many class fields is finite.
\end{proof}

An explicit discriminant estimate would give a quantitative route
to the hypothesis. For example, with \(\mathfrak d_{K/F_0}\)
the relative discriminant ideal, an inequality
\begin{equation}\label{eq:discriminant-target}
 \log N_{F_0/\Q}\mathfrak d_{F_0(\xi)/F_0}
       \le A_D q(\xi)+B_D
\end{equation}
on integral classes of degree at most \(D\) would suffice.
The qualitative ramification condition only asks for some finite
set depending on \(B,D\); it is weaker than
\MathEqref{eq:discriminant-target}.

\subsection{Why uniform fixed-field estimates do not suffice}
\label{sec:quadratic-example}

\begin{example}\label[example]{ex:quadratic-fields}
Choose infinitely many distinct quadratic extensions \(K_i/F_0\)
with characters \(\chi_i\). Put
\[
 \mathcal G=\bigoplus_i\Q e_i,\qquad
 \sigma e_i=\chi_i(\sigma)e_i,\qquad
 \Lambda(E)=\bigoplus_{K_i\subseteq E}\Z e_i,
 \qquad q\left(\sum a_ie_i\right)=\sum a_i^2.
\]
The lattices are compatible with restriction and trace. Every
finite extension \(E\) has at most \([E:F_0]-1\) distinct
quadratic subfields: their compositum is multiquadratic, and the
number of its quadratic subfields is one less than its degree.
Consequently
\[
 \rk\Lambda(E)\le[E:F_0]-1,\qquad
 q(\lambda)\ge1\ (\lambda\ne0),\qquad
 |\Lambda(E)/m\Lambda(E)|\le m^{[E:F_0]-1}.
\]
All fixed-field height balls are finite. Nevertheless the classes
\(e_i\in\Lambda(K_i)\) all have degree two and height one.
Thus bounded-degree Northcott fails despite uniform
degree-dependent bounds on ranks, gaps and Kummer quotients.
\end{example}

The fixed-\(S\) cohomological statement of
\Cref{prop:fixed-S} holds over a chosen field of definition.
\Cref{thm:ramification} instead bounds the fields themselves.
This distinction is essential when using residual arithmetic
regulators to study the conjectural Northcott property.
